\documentclass[APA,STIX1COL]{WileyNJD-v2}

\articletype{Special Issue Paper}

\received{}
\revised{}
\accepted{}

\raggedbottom

\begin{document}

\title{Extreme Maximum Temperature Events in Mexico: A Weibull Distribution Approach with Optimized Plotting Positions}

\author[1]{Sazcha Olivera}

\authormark{OLIVERA \textsc{et al}}

\address[1]{\orgdiv{Departamento de Teoría y Procesos del Diseño, Ciencias de la Comunicación y Diseño}, \orgname{Universidad Autónoma Metropolitana-Unidad Cuajimalpa}, \orgaddress{\state{Ciudad de México}, \country{México}}}

\corres{Sazcha Olivera, Departamento de Teoría y Procesos del Diseño, Ciencias de la
Comunicación y Diseño, Universidad Autónoma Metropolitana-Unidad Cuajimalpa, 05370 Ciudad de México, Mexico. \email{satzcha@msn.com}}

%% ─────────────────────────────────────────────────────────────
\abstract[Resumen]{
Se propone una metodología adaptativa para modelar la cola derecha de excedencias de temperatura máxima diaria mediante la distribución de Weibull truncada de dos parámetros, extendiendo el esquema ponderado de \citet{olivera2019} en tres dimensiones:
(i)~el umbral de referencia se estima como la temperatura modal mediante estimación kernel de la densidad (KDE); (ii)~el número de clases del histograma agrupado se determina por la regla de Freedman--Diaconis; y (iii)~el parámetro~$a$ de la familia simétrica $p_i(a)=(i-a)/(n+1-2a)$ se optimiza
adaptando la posición de trazado a los datos, operación que se justifica teóricamente para la subfamilia $0<k<1$. La metodología se aplica a temperaturas máximas diarias de cinco ciudades de México y el suroeste de Estados Unidos en tres periodos de veinte años; los resultados muestran que el umbral modal~$\mu^*$ presenta una tendencia
ascendente en la mayoría de las ciudades analizadas, mientras que el
parámetro de forma~$\hat\beta$ permanece comparativamente estable dentro
de cada región, evidenciando que los cambios en los extremos térmicos se
expresan principalmente como desplazamientos del nivel habitual de temperatura.
}

\keywords{distribuci\'on Weibull; posiciones de trazado;
par\'ametro~$a$; estimaci\'on kernel de la densidad; temperaturas
extremas; cola Weibull; construcci\'on de clases;
Freedman--Diaconis; pesos de Bergman.}

\jnlcitation{\cname{\author{Olivera S.}
             (\cyear{2026}),
             \ctitle{Extreme Maximum Temperature Events in Mexico: A Weibull Distribution Approach with Optimized Plotting Positions},
             \cjournal{[Revista]}, \cvol{00:1--XX}.}

\maketitle

%% ════════════════════════════════════════════════════════════
\section{Introducción}\label{sec:intro}
%% ════════════════════════════════════════════════════════════

La distribución de Weibull ocupa una posición particular en la estadística aplicada porque combina una formulación relativamente
simple con una notable flexibilidad de forma. Desde sus primeras aplicaciones al análisis de fallas de materiales \citep{weibull1951}, su uso se ha extendido a la hidrología, la
climatolog\'ia, la ingenier\'ia de confiabilidad y el an\'alisis de
riesgos. En todos esos campos, el inter\'es no recae en el centro
de la distribuci\'on sino en su cola derecha: la probabilidad de
que la variable supere umbrales elevados, raramente observados pero
con consecuencias significativas.

Cuando la distribuci\'on de Weibull se estima mediante el gr\'afico
de probabilidad, las observaciones ordenadas no aportan, por s\'i
mismas, sus probabilidades acumuladas. A cada $x_{(i)}$ debe
asign\'arsele una probabilidad de no excedencia, denominada
\emph{posici\'on de trazado}~$p_i$. La elecci\'on de $p_i$ no es
inocua: es la que determina las coordenadas de la regresi\'on lineal
en el papel de probabilidad y, por tanto, las estimaciones de los
par\'ametros de forma~$k$ y escala~$\lambda$.

Una familia que re\'une muchas de las reglas cl\'asicas es
\begin{equation}\label{eq:familia_intro}
p_i(a) = \frac{i - a}{n + 1 - 2a},\qquad 0 \leq a < \tfrac{1}{2}.
\end{equation}
El valor $a=0$ genera la regla de Weibull~\citep{weibull1951};
$a\approx 0.30$ corresponde a la aproximaci\'on de Benard y
Bos-Levenbach~\citeyearpar{benard1953}; $a=0.375$ conduce a la
f\'ormula de Blom~\citeyearpar{blom1958}; $a\approx 0.44$ a
Gringorten~\citeyearpar{gringorten1963}, y el l\'imite $a\to 0.5$
reproduce la posici\'on de Hazen. Aunque estas reglas difieren poco
num\'ericamente, su impacto sobre los cuantiles extremos estimados resulta ser sustancial, particularmente cuando el tama\~no de muestra
es moderado y la distribuci\'on tiene una cola pesada.

Existe, al mismo tiempo, un fundamento te\'orico bien establecido
para las posiciones de trazado. Si $X_{(i)}$ es el i-ésimo estadístico de orden de una muestra continua, entonces definimos $U_{(i)}=F(X_{(i)})$ tal que 

\begin{equation*}
    U_{(i)}\sim\mathrm{Beta}(i,\,n+1-i)
\end{equation*}

Y 

\begin{equation}
    \mathbb{E}[U_{(i)}]=\dfrac{i}{n+1}, 
\end{equation}

\medskip

lo que explica la regla de Weibull
\citep{makkonen2014}. Sin embargo, el ajuste linealizado utiliza la
transformaci\'on doble logar\'itmica $g(p)=\ln[-\ln(1-p)]$, que es
estrictamente convexa. En consecuencia, la desigualdad de Jensen
implica que $g(E[U_{(i)}])\neq E[g(U_{(j)})]$, y la posici\'on
que es exacta en la escala de probabilidades no minimiza el sesgo
en la escala transformada. La magnitud de esta brecha depende de
la curvatura de~$g$ en el rango de probabilidades relevante, y
por tanto del par\'ametro de forma~$k$ del modelo ajustado.

\citet{goda1988,goda2010} determinaron mediante simulaci\'on Monte
Carlo que la posici\'on \'optima para extremos Weibull depende
expl\'icitamente de~$k$, proporcin\'ando coeficientes calibrados
para el an\'alisis de oleaje extremo. Kim et~al.
\citeyearpar{kim2012,kim2015} extendieron ideas similares a la
distribuci\'on log\'istica generalizada. A pesar de estos
precedentes, la pr\'actica habitual en climatolog\'ia e hidrolog\'ia
contin\'ua fijando~$a$ de antemano o seleccionando la posici\'on
que produce el mayor~$R^2$ dentro de la misma muestra que se analiza,
lo que introduce optimismo de selecci\'on sin garant\'ia de
generalizaci\'on.

El punto de partida inmediato de esta investigaci\'on es
\citet{olivera2019}. Su procedimiento compar\'o posiciones de
trazado para el an\'alisis de excedencias de precipitaci\'on y
utiliz\'o la funci\'on de peso derivada de Bergman
\citeyearpar{bergman1986} para ponderar el ajuste linealizado:

\begin{equation}\label{eq:bergman_intro}
w_i = (1 - p_i)\,[\ln(1-p_i)]^2.
\end{equation}

\medskip

Esta precisión tambi\'en revela una caracter\'istica del peso que
resulta pertinente para el an\'alisis de extremos: $w_i$ alcanza su
m\'aximo en $p\approx 0.8647$ y tiende a cero cuando $p\to 1$, por
lo que no asigna el mayor peso al m\'aximo de la muestra. El trabajo
de \citet{olivera2019} mostr\'o que la elecci\'on de la posici\'on
de trazado puede cambiar las conclusiones sobre tendencias en
cuantiles extremos; sin embargo, trat\'o~$a$ como fijo en todas las
muestras.

La propuesta de la presente investigaci\'on consiste en estimar una
regla $\hat{a}=g(n,\tilde{k})$ dentro de la familia
sim\'etrica~\eqref{eq:familia_intro}, donde $n$ es el n\'umero de
clases y $\tilde{k}$ es una estimaci\'on preliminar del par\'ametro
de forma obtenida sin posiciones de trazado. El dominio principal
es la Weibull con $0<k<1$: en este intervalo la distribuci\'on
presenta una cola exponencial estirada que decrece m\'as lentamente
que una exponencial, aunque no constituye una cola de potencia
(Pareto). Este r\'egimen de cola Weibull \citep{gardes2005,
diebolt2008} es el emp\'iricamente relevante para excedencias de
temperatura m\'axima en climas continentales y semi\'aridos, donde
los episodios de calor intenso son frecuentes pero el calor extremo
tiene un l\'imite f\'isico superior que descarta las leyes de potencia.

%%NUEVA METODOLOGÍA

La metodología integra, además de la optimización de~$a$, tres contribuciones sobre el preprocesamiento de los datos que responden a limitaciones de las implementaciones anteriores:

\textbf{1. Umbral modal mediante KDE.}
El umbral que separa días normales de
excedencias se estima como la temperatura modal mediante el estimador kernel de la densidad~\citep{rosenblatt1956,parzen1962}. Esto permite reflejar la temperatura más
frecuente y hace posible que la proporción de excedencias sea una cantidad estimada, no fijada de antemano.

\textbf{2. Número de clases por Freedman--Diaconis.}
El ancho del histograma empleado en la regresión agrupada se determina mediante la regla de Freedman-Diaconis~\citeyearpar{freedman1981},
que es asintóticamente óptima bajo condiciones más débiles que la regla de Sturges y adapta el número de clases~$n$ al
tamaño y la dispersión de los datos de cada ciudad y periodo.

\textbf{3. Modelo de mezcla.}
La variable de excedencia $T_i=\max(X_i-\mu^*,0)$ no es puramente continua: tiene una masa puntual en $T=0$ para todos los días con
$X_i\leq\mu^*$ y una componente Weibull continua para $T>0$. La especificación explícita de la distribución de mezcla

\begin{equation}\label{eq:mezcla_intro}
F_{\mathrm{mix}}(t) = p_0 + (1-p_0)\,F(t;\,k,\lambda),\qquad t>0,
\end{equation}

es necesaria para que la función de distribución empírica y el modelo ajustado se comparen en la misma escala. Sin este paso,
la FDA empírica parte de~$p_0\approx 0.5$ mientras que la FDA Weibull parte de~0, produciendo residuos artificialmente grandes.

La pregunta de investigación central es: \textit{\textquestiondown Es posible cuantificar cambios en la cola derecha de la distribución de los eventos de temperatura máxima a lo largo de varias décadas, utilizando la distribución de Weibull ajustada a las excedencias?} A partir de esta pregunta se plantean dos hipótesis verificables:

\begin{enumerate}
\item[\textbf{H1.}] \textbf{Estabilidad de la forma distribucional.}
  El parámetro de forma~$k$ de la distribución Weibull ajustada a las excedencias de temperatura máxima no varía significativamente entre los tres periodos analizados dentro de cada ciudad. Los cambios en los extremos térmicos se reflejan principalmente en desplazamientos del umbral modal~$\mu^*$ y de la escala~$\lambda$, no en cambios en la forma de la distribución.

\item[\textbf{H2.}] \textbf{Desplazamiento de los niveles extremos.}
  El umbral modal~$\mu^*$ exhibe una tendencia positiva entre periodos
  en la mayor\'ia de las ciudades analizadas. Como consecuencia, aun
  cuando la forma de la distribuci\'on permanezca estable, las
  probabilidades de superar umbrales cr\'iticos fijos aumentan entre
  d\'ecadas, con repercusiones sobre la demanda de energ\'ia para
  climatizaci\'on en entornos urbanos.
\end{enumerate}

La metodolog\'ia propuesta permite contrastar ambas hip\'otesis mediante
la comparaci\'on de los par\'ametros estimados~$(k, \lambda, \mu^*)$
y de las probabilidades de excedencia calculadas para umbrales
cl\'imaticamente relevantes.

%% ════════════════════════════════════════════════════════════
\section{Distribución de Weibull y el régimen de cola}\label{sec:weibull}
%% ════════════════════════════════════════════════════════════

Para evitar la ambigüedad frecuente en la literatura sobre el
intercambio de los símbolos~$\alpha$ y~$\beta$ entre tradiciones
distintas, se adopta la notación internacional: $k>0$ para el
parámetro de forma y $\lambda>0$ para el de escala
\citep[véase también][]{makkonen2014}. La función de distribución
acumulada (FDA) y la función de supervivencia de la distribución
Weibull de dos parámetros son
\begin{equation}\label{eq:fda}
F(x;\,k,\lambda)
  = 1 - \exp\!\left\{-\left(\frac{x}{\lambda}\right)^{\!k}\right\},
  \qquad x > 0,
\end{equation}
\begin{equation}\label{eq:supervivencia}
S(x;\,k,\lambda)
  = \exp\!\left\{-\left(\frac{x}{\lambda}\right)^{\!k}\right\},
\end{equation}
con tasa de riesgo
\begin{equation}\label{eq:hazard}
h(x) = \frac{k}{\lambda}\left(\frac{x}{\lambda}\right)^{k-1}.
\end{equation}

\medskip
Esta función es decreciente para $0<k<1$, constante para $k=1$
(distribución exponencial) y creciente para $k>1$.

Siguiendo \citet{gardes2005} y \citet{diebolt2008}, se dice que una
distribución pertenece a la \emph{clase de cola Weibull} si su función
de supervivencia satisface
\begin{equation}\label{eq:cola_weibull}
-\ln S(x) = x^\tau\,\ell(x),\qquad \tau>0,
\end{equation}

\medskip

donde $\ell$ es una función de variación lenta en infinito. La Weibull
estándar es el miembro canónico de esta clase con $\tau=k$ y
$\ell\equiv\lambda^{-k}$. Para $0<k<1$, la cola es \emph{más pesada
que la exponencial} pero \emph{más ligera que cualquier ley de
potencia} (tipo Pareto); este régimen intermedio descarta tanto las
distribuciones de cola pesada, como la distribución
exponencial pura ($k=1$). En $k>1$, por su parte, se describen colas que disminuyen rápidamente; bajo excedencias de temperatura sobre el umbral modal $\mu^*$, puede reflejar
que eventos muy grandes son infrecuentes una vez que se supera la temperatura típica.

El régimen de~$k$ tiene además una implicación directa sobre la
elección de la posición de trazado. \citet{goda1988,goda2010}
mostraron mediante simulación Monte Carlo que la posición óptima para
extremos Weibull depende de~$k$ y que la dependencia es más pronunciada
cuando $k<1$, lo que motiva el enfoque adaptativo propuesto.

%% ════════════════════════════════════════════════════════════
\section{Umbral modal y modelo de mezcla}\label{sec:threshold}
%% ════════════════════════════════════════════════════════════

\subsection{Motivación}

Sea $\{X_i\}_{i=1}^{N}$ la serie de temperaturas m\'aximas diarias
observadas en una ciudad y per\'iodo determinados. Para construir
la variable de excedencia es necesario fijar un valor de referencia
$\mu^*$ que represente la temperatura ``habitual'': los d\'ias con
$X_i\leq\mu^*$ se clasificar\'an como normales y los d\'ias con
$X_i>\mu^*$ como excedencias. La elecci\'on de $\mu^*$ condiciona qu\'e proporci\'on
de las observaciones contribuye a la estimaci\'on de la cola y, por
tanto, los par\'ametros de la distribuci\'on ajustada.

La alternativa más utilizada en la literatura es la mediana
muestral~$\hat{m}$. Su limitación principal es que, por construcción, impone que aproximadamente el 50\% de los días sean excedencias,
independientemente de la asimetría real de la distribución. En
ciudades con cola derecha extendida ---como Chihuahua o Phoenix, donde
los episodios de calor extremo desplazan el percentil~50 hacia
arriba--- esto provoca que días de calor moderado queden
clasificados como excedencias, deformando los eventos genuinamente atípicos o extremos.

Se propone utilizar en cambio la \emph{moda} de la distribución como un
umbral~$\mu^*$ que refleja el valor de temperatura que ocurre con mayor
frecuencia. La moda no impone a priori ninguna fracción de
excedencias; el valor $\hat{p}_0 = P(T=0)$ resulta de los propios
datos. Sin embargo, la moda empírica directa es inestable, pues pequeños
cambios en el criterio de redondeo pueden desplazarla varios décimos de
grado, y con miles de registros puede no ser única. El histograma de
frecuencias evita el problema del redondeo pero a su vez introduce arbitrariedad
en la elección de los límites de los intervalos y produce una curva
escalonada, inadecuada para localizar un máximo continuo.

La \emph{estimación kernel de la densidad} (KDE) resuelve ambas
dificultades: suaviza la distribución empírica, produciendo una curva diferenciable cuyo máximo constituye el estimador $\mu^*$ propuesto.

\subsection{El estimador kernel de la densidad}\label{subsec:kde}

Sea $K:\mathbb{R}\to[0,\infty)$ una funci\'on kernel y $h>0$ el ancho de banda que
controla qu\'e tan amplia es cada campana. El KDE en un punto
gen\'erico~$x$ se define como \citep{rosenblatt1956,parzen1962}:

\begin{equation}\label{eq:kde}
\hat{f}_h(x)
  = \frac{1}{Nh}\sum_{i=1}^{N}
    K\!\left(\frac{x - X_i}{h}\right).
\end{equation}

\medskip

Se adopta el n\'ucleo gaussiano est\'andar:
\begin{equation}\label{eq:kde_gauss}
K(u) = \frac{1}{\sqrt{2\pi}}\,e^{-u^2/2},
\end{equation}

\medskip

con lo que cada observaci\'on contribuye una densidad normal con
media $X_i$ y desviaci\'on est\'andar $h$. El estimador es
entonces la mezcla equiponderada de $N$ distribuciones normales,
una centrada en cada observaci\'on.

La elecci\'on del n\'ucleo gaussiano es est\'andar en la pr\'actica, pues produce curvas
diferenciables en todo punto y no impone un soporte compacto,
lo que la hace computacionalmente m\'as conveniente y
estad\'isticamente preferible cuando la densidad subyacente
es continua y sin discontinuidades abruptas como la temperatura.

\subsubsection{Regla de Scott y estimación del umbral modal}
El par\'ametro~$h$ determina la suavidad de la curva. Un valor
de~$h$ muy peque\~no reproduce cada observaci\'on individualmente mientras que un~$h$
muy grande aplana la curva hasta perder la estructura de la
distribuci\'on. Existe un valor \'optimo que balancea el sesgo y la varianza.

Formalmente, el error cuadr\'atico integrado medio tiene la
descomposici\'on asint\'otica

\begin{equation}\label{eq:mise}
\mathrm{MISE}(h)
  \approx \underbrace{\frac{h^4}{4}
            \int[f''(x)]^2\,dx}_{\mathrm{sesgo}^2}
  + \underbrace{\frac{1}{Nh}
            \int K^2(u)\,du}_{\mathrm{varianza}}.
\end{equation}

\medskip

Minimizando respecto a~$h$ se obtiene que el \'optimo satisface
$h^*\propto N^{-1/5}$, un resultado universal que no depende del
n\'ucleo ni de la densidad desconocida. \citet{scott1992} evaluó
la integral $\int[f''(x)]^2\,dx$ bajo el supuesto de densidad
gaussiana subyacente y n\'ucleo gaussiano, obteniendo

\begin{equation}\label{eq:scott}
h_{\mathrm{Scott}} = 1.059\,\hat{\sigma}\,N^{-1/5},
\end{equation}

\medskip

donde $\hat{\sigma}$ es la desviaci\'on est\'andar muestral.

La regla de Scott es la opci\'on est\'andar cuando la densidad subyacente es aproximadamente unimodal
\citep{silverman1986}; sus supuestos son aceptables para la
distribuci\'on de temperatura m\'axima en los per\'iodos
analizados.

\medskip

Con el estimador $\hat{f}_h$ disponible, la moda de la distribuci\'on
se estima como el argumento que maximiza la densidad suavizada:

\begin{equation}\label{eq:mustar}
\mu^* = \underset{x\in\mathcal{G}}{\arg\max}\;\hat{f}_h(x),
\end{equation}

donde $\mathcal{G}$ es una rejilla uniforme de $J=2{,}000$ puntos
en $[\min X_i, \max X_i]$. La notaci\'on $\arg\max$ denota el punto $x$ donde $\hat{f}_h$ alcanza su valor m\'aximo. Con $J=2{,}000$ puntos
en un rango t\'ipico de $40\,{}^\circ\mathrm{C}$, la resoluci\'on
es $\Delta x=0.02\,{}^\circ\mathrm{C}$, cinco veces m\'as fina
que la precisi\'on de los term\'ometros, de modo que los errores
de discretizaci\'on son despreciables.

\subsection{La variable de excedencia}\label{subsec:excedencia}

A partir de $\mu^*$ se define la variable de excedencia

\begin{equation}\label{eq:Ti}
T_i = \max(X_i - \mu^*,\;0),\qquad i=1,\ldots,N.
\end{equation}

\medskip

La variable $T_i$ mide cu\'antos grados supera la temperatura del
d\'ia~$i$ al valor modal de referencia. Para los d\'ias con
$X_i\leq\mu^*$, se tiene $T_i=0$. Para los d\'ias con
$X_i>\mu^*$, la excedencia es la diferencia positiva
$T_i=X_i-\mu^*$, que puede tomar cualquier valor continuo en
$(0,\infty)$. Esta mezcla de un valor puntual y una componente
continua requiere un modelo distribucional especial.

Así pues, cuando $T>0$, la
distribuci\'on condicional de~$T$ dado~$T>0$ se modela como Weibull
con par\'ametros~$(k,\lambda)$:

\begin{equation}\label{eq:dist_cond}
P(T\leq t \mid T>0) = F(t;\,k,\lambda),\qquad t>0.
\end{equation}

El par\'ametro~$k$ en la ecuaci\'on~\eqref{eq:dist_cond} caracteriza
la forma de la cola derecha de la distribuci\'on de temperatura
original por encima de~$\mu^*$. Estimar~$k$ y~$\lambda$ con
precisi\'on a partir de datos agrupados es el problema inferencial
central que aborda la metodolog\'ia de la
Secci\'on~\ref{sec:estimation}.


\subsection{El modelo de mezcla}

Las distribuciones de Weibull est\'andar se definen sobre
$(0,\infty)$ y son absolutamente continuas: asignan probabilidad
cero a cualquier valor puntual aislado. Si se intenta ajustar
directamente una Weibull a la variable~$T_i$ sin tratar la masa
en cero, se produce una inconsistencia, pues la distribuci\'on
te\'orica dice que $P(T=0)=0$, mientras que los datos muestran
$P(T=0)\approx 0.5$. El m\'etodo de ajuste intenta compensar esta discrepancia deformando los par\'ametros de la distribuci\'on,
produciendo estimaciones sesgadas de~$k$ y~$\lambda$.

Tomando la ley de probabilidad total como punto de partida, para cualquier valor $t>0$, la probabilidad de que $T\leq t$
puede descomponerse condicionando en los dos eventos mutuamente
excluyentes $\{T=0\}$ y $\{T>0\}$:

\begin{align}
P(T\leq t)
  &= P(T\leq t\mid T=0)\,P(T=0)
   + P(T\leq t\mid T>0)\,P(T>0) \notag \\
  &= 1\cdot p_0 + P(T\leq t\mid T>0)\cdot (1-p_0),
    \label{eq:pt_total}
\end{align}

\medskip

donde el primer t\'ermino es~1 porque si $T=0$ y $t>0$, se
cumple siempre $T\leq t$, y se ha denotado $p_0=P(T=0)$.
La distribuci\'on condicional $P(T\leq t\mid T>0)$ es lo que
se modela mediante la Weibull de dos par\'ametros:
$P(T\leq t\mid T>0)=F(t;\,k,\lambda)$.
Sustituyendo en~\eqref{eq:pt_total} se obtiene el
\emph{modelo de mezcla}:

\begin{equation}\label{eq:mezcla}
F_{\mathrm{mix}}(t;\,k,\lambda)
  = \begin{cases}
    p_0 & t=0,\\[6pt]
    p_0 + (1-p_0)\,F(t;\,k,\lambda) & t>0.
    \end{cases}
\end{equation}

\medskip

Los par\'ametros de la Weibull se interpretan ahora
exclusivamente sobre los d\'ias que son excedencias positivas:
$k$ caracteriza c\'omo decrece la probabilidad de superar un
umbral dado dentro de la poblaci\'on de d\'ias calurosos,
y $\lambda$ establece la escala t\'ipica de esas excedencias.

\medskip

La probabilidad de masa en cero se estima directamente de
los datos:

\begin{equation}\label{eq:p0}
\hat{p}_0
  = \frac{\#\{i : X_i\leq\mu^*\}}{N}.
\end{equation}

\medskip

A diferencia del caso con la mediana como umbral, donde
$\hat{p}_0=0.5$ exactamente por construcci\'on, con el umbral
modal $\hat{p}_0$ es una cantidad estimada que depende de la forma real de la distribuci\'on. En ciudades con
distribuciones menos asim\'etricas, $\hat{p}_0$ puede acercarse
a 0.5; en climas con cola calurosa pronunciada, puede alejarse
sustancialmente como se observa en la Tabla \ref{tab:p0}.

\begin{table}[ht]
\centering
\begin{tabular}{lccc}
\toprule
Ciudad       & 1940--1968 & 1969--1996 & 1997--2026 \\
\midrule
CDMX         & $\hat{p}_0$ = 0.4796  & $\hat{p}_0$ = 0.4979 & $\hat{p}_0$ = 0.4782 \\
Chihuahua    & $\hat{p}_0$ = 0.7612  & $\hat{p}_0$ = 0.6624 & $\hat{p}_0$ = 0.6828 \\
Guadalajara  & $\hat{p}_0$ = 0.4066 & $\hat{p}_0$ = 0.4321 & $\hat{p}_0$ = 0.4693 \\
Monterrey    & $\hat{p}_0$ = 0.7506 & $\hat{p}_0$ = 0.6372 & $\hat{p}_0$ = 0.6759 \\
Phoenix      & $\hat{p}_0$ = 0.8199 & $\hat{p}_0$ = 0.8180 & $\hat{p}_0$ = 0.8190 \\
\bottomrule
\end{tabular}
\caption{Probabilidad empírica de la clase cero $\hat{p}_0 = P(T=0)$
por ciudad y período.}
\label{tab:p0}
\end{table}.

%% ════════════════════════════════════════════════════════════
\section{Construcción de clases y el papel de~$n$}\label{sec:classes}
%% ════════════════════════════════════════════════════════════

La regresi\'on WLS de la Secci\'on~\ref{sec:estimation} opera sobre
datos agrupados: las excedencias positivas $\{T_i:T_i>0\}$ se
dividen en $n$~clases, cada una caracterizada por un valor
representativo~$t_j$ y una frecuencia~$f_j$. El n\'umero de
clases~$n$ es un par\'ametro impl\'icito del procedimiento
completo: afecta las posiciones de trazado asignadas a cada clase,
los pesos empleados en la regresi\'on y, en \'ultimo t\'ermino,
los par\'ametros~$(k,\lambda)$ estimados y la bondad de ajuste.

En el contexto de temperatura, con datos que abarcan varias d\'ecadas y cinco
ciudades con distribuciones distintas, un~$n$ fijo universal
ser\'ia inapropiado. Se le trata por tanto como una cantidad
determinada por los datos.

Se distinguen dos tipos de clase.

La primera clase, que contiene todas las observaciones con $T_i = 0$, corresponde a la masa de probabilidad acumulada por todos los
d\'ias en que la temperatura no supera el umbral modal $\mu^*$. Su
peso emp\'irico es $\hat{p}_0 = P(T = 0) = P(X \leq \mu^*)$. Su valor representativo se fija en
$t_1=\varepsilon=0.01\,{}^\circ\mathrm{C}$ para permitir la transformaci\'on logar\'itmica requerida por la linealizaci\'on
Weibull.

Esta clase se introduce para modelar expl\'icitamente la estructura
mixta de la variable $T = \max(X - \mu^*, 0)$, ya que la distribuci\'on
de mezcla $F_{\text{mix}}$ requiere que la masa en cero est\'e separada de la componente
Weibull continua. La clase cero no participa en la regresi\'on lineal
ponderada (WLR) que estima los par\'ametros $k$ y $\lambda$; su \'unica
funci\'on es asegurar que la FDA emp\'irica y la FDA ajustada se comparen
en la misma escala.

\medskip

El segundo tipo de clases, las positivas, son los intervalos semibiertos
$(0,h_{\mathrm{FD}}]$, $(h_{\mathrm{FD}},2h_{\mathrm{FD}}]$,
$\ldots$, con valores representativos $t_2=h_{\mathrm{FD}}$,
$t_3=2h_{\mathrm{FD}}$, etc.

\medskip

El ancho de clase $h_{\mathrm{FD}}$ se determina mediante la regla
de Freedman--Diaconis \citep{freedman1981}:

\begin{equation}\label{eq:FD}
h_{\mathrm{FD}}
  = \frac{2\,\mathrm{IQR}(T)}{N_+^{1/3}},
\end{equation}

\medskip

donde $\mathrm{IQR}(T)=Q_{0.75}(T)-Q_{0.25}(T)$ es el rango
intercuart\'ilico de todas las excedencias $T_i\geq 0$ y
$N_+$~es el n\'umero de excedencias positivas.

\medskip

Partiendo de la clase cero y las clases positivas generadas por
la ecuaci\'on~\eqref{eq:FD}, la construcci\'on contin\'ua hasta
encontrar clases vac\'ias. Asimismo, las clases con frecuencia nula se eliminan de la cola para obtener
el conjunto activo de $n$~clases. La \'ultima clase ---la que lleva
la frecuencia acumulada a 1--- tambi\'en se excluye de
la regresi\'on WLS, pues requerir\'ia evaluar
$Y_j=\ln[-\ln(1-p_j)]$ en $p_j=1$, que es indefinido.

%% ════════════════════════════════════════════════════════════
\section{Estimaci\'on WLS y optimizaci\'on
         adaptativa de~$a$}\label{sec:estimation}
%% ════════════════════════════════════════════════════════════

La FDA Weibull permite una linearización exacta. Tomando logaritmos dos veces en la ecuaci\'on~\eqref{eq:fda}:

\begin{align}
1 - F(t) &= \exp\!\left\{-\!\left(\frac{t}{\lambda}\right)^{\!k}\right\}
  \nonumber\\[4pt]
-\ln[1-F(t)] &= \left(\frac{t}{\lambda}\right)^{\!k}
  \nonumber\\[4pt]
\ln\!\bigl[-\ln(1-F(t))\bigr] &= k\ln t - k\ln\lambda.
  \label{eq:doble_log}
\end{align}

\medskip

Definiendo las variables transformadas para la $j$-\'esima clase positiva
con valor representativo~$t_j$ y posici\'on de trazado~$p_j$:

\begin{equation}\label{eq:XY}
X_j = \ln(t_j),\qquad Y_j = \ln[-\ln(1-p_j)],
\end{equation}

\medskip
la ecuaci\'on~\eqref{eq:doble_log} queda como una regresi\'on lineal exacta:
\begin{equation}\label{eq:lineal}
Y_j = k\,X_j + b,\qquad b = -k\ln\lambda.
\end{equation}

\medskip

La pendiente de la regresi\'on~\eqref{eq:lineal} estima directamente~$k$,
y el intercepto~$b$ permite recuperar~$\lambda$:

\begin{equation}\label{eq:parametros}
\hat{k} = \frac{S_{XY}}{S_{XX}},\qquad
\hat{\lambda} = \exp\!\left(-\frac{b}{k}\right)
              = \exp\!\left(-\frac{\bar{Y}_w - \hat{k}\,\bar{X}_w}
                                   {\hat{k}}\right),
\end{equation}

\medskip

donde las medias ponderadas se definen como

\begin{align}
\bar{X}_w &= \frac{\sum_j w_j X_j}{\sum_j w_j},\qquad
\bar{Y}_w  = \frac{\sum_j w_j Y_j}{\sum_j w_j}
\end{align}

\medskip

y las sumas de productos son

\begin{align}
S_{XY} &= \sum_j w_j(X_j - \bar{X}_w)(Y_j - \bar{Y}_w),\quad
S_{XX}  = \sum_j w_j(X_j - \bar{X}_w)^2.\label{eq:sumas_w}
\end{align}

\medskip
El coeficiente de determinaci\'on de la regresi\'on es

\begin{equation}\label{eq:R2_wlr}
R^2_{\mathrm{WLR}} = \frac{S_{XY}^2}{S_{XX}\,S_{YY}},\qquad
S_{YY} = \sum_j w_j(Y_j - \bar{Y}_w)^2.
\end{equation}

\medskip

Finalmente, siguiendo \citet{olivera2019} y \citet{bergman1986}, se considera la
función de peso:

\begin{equation}\label{eq:peso_A}
w_j = (1 - p_j)\,[\ln(1-p_j)]^2,
\end{equation}

\subsection{Estimación adaptativa de \textbf{$a$}}\label{subsec:familia}

La familia sim\'etrica de un par\'ametro~$a\in[0,\tfrac{1}{2})$ engloba
los estimadores cl\'asicos:
\begin{equation}\label{eq:familia}
p_i(a) = \frac{i-a}{n + 1 - 2a},\qquad i=1,\ldots,n.
\end{equation}

\medskip
La posici\'on de la clase cero se fija siempre en
$p_1^{(0)}=1/(n+1)$, correspondiente a $a=0$ (regla de Weibull).
Esto asegura que la elecci\'on de~$a$ para las clases positivas no
propague un sesgo a trav\'es de la clase que queda excluida de la regresi\'on.

Para las clases positivas, el parámetro~$a$ se determina mediante un
procedimiento de optimización numérica, seleccionando el valor que
maximiza la bondad de ajuste del modelo Weibull. El parámetro obtenido
se emplea posteriormente en la estimación de los parámetros de forma y
escala de la distribución.

\subsection{Algoritmo operacional}\label{subsec:algoritmo}

El procedimiento propuesto para el cálculo de los parámetros Weibull $k$ y $\lambda$ se resume a continuación.

\begin{enumerate}
\item Estimar $\mu^*$ mediante KDE (ecuaciones~\eqref{eq:kde}--\eqref{eq:mustar})
  y calcular $T_i=\max(X_i-\mu^*,0)$.
\item Determinar $h_{\mathrm{FD}}$ (ecuaci\'on~\eqref{eq:FD}) y
  construir la clase cero y las clases positivas.
\item Eliminar las clases vac\'ias de la cola e identificar $n$.
\item Obtener estimaciones preliminares $(\tilde{k},\tilde{\lambda})$ mediante m\'axima verosimilitud.
\item Refinar los par\'ametros mediante optimizaci\'on conjunta, obteniendo
$(k,\lambda)$.
\end{enumerate}

%% ════════════════════════════════════════════════════════════
\section{\'Area de estudio y datos}
%% ════════════════════════════════════════════════════════════

Se obtuvieron registros de temperatura máxima diaria para cinco ciudades con climas distintos dentro de América del
Norte: Ciudad de México~(CDMX), Guadalajara, Monterrey, Chihuahua y Phoenix~(AZ, EUA). Las cuatro primeras se ubican en diferentes
elevaciones y latitudes dentro de México; Phoenix representa un clima desértico árido a baja elevación en el suroeste de los
Estados Unidos.

Los datos de temperatura m\'axima diaria se obtuvieron del conjunto
de datos \emph{ERA5 hourly data on single levels
from 1940 to present} \citep{hersbach2023era5}, conjunto de
rean\'alisis clim\'atico global producido por el Centro Europeo de
Predicci\'on Meteorol\'ogica a Plazo Medio (ECMWF) y distribuido a
trav\'es del Climate Data Store del Servicio de Cambio Clim\'atico
de Copernicus~(C3S) \citep{hersbach2020era5}, y comprenden observaciones diarias de
temperatura máxima para los tres períodos analizados:
1940--1968, 1969--1996 y 1997--2026. Las cinco ciudades disponen
de registros completos. El número de
observaciones por ciudad y período varía dependiendo del número de días por año.

La Figura~\ref{fig:kde} muestra la densidad KDE estimada de la
temperatura máxima diaria para cada ciudad y período, junto con
el umbral modal~$\mu^*$ identificado en cada caso. Se observa que
la distribución de temperatura varía considerablemente entre
ciudades y que,
dentro de cada ciudad, la posición del pico de la densidad (umbral
modal) muestra una tendencia ascendente entre períodos en la mayoría
de los casos.

\begin{figure}[ht]
    \centering

    % Fila 1
    \begin{minipage}{0.33\textwidth}
        \centering
        \includegraphics[width=\textwidth]{images/fig1_kde_cdmx_v2.png}
    \end{minipage}
    \hfill
    \begin{minipage}{0.33\textwidth}
        \centering
        \includegraphics[width=\textwidth]{images/fig1_kde_chihuahua_v2.png}
    \end{minipage}
    \hfill
    \begin{minipage}{0.33\textwidth}
        \centering
        \includegraphics[width=\textwidth]{images/fig1_kde_guadalajara_v2.png}
    \end{minipage}

    \vspace{0.5cm}

    % Fila 2
    \begin{minipage}{0.33\textwidth}
        \centering
        \includegraphics[width=\textwidth]{images/fig1_kde_monterrey_v2.png}
    \end{minipage}
    \begin{minipage}{0.33\textwidth}
        \centering
        \includegraphics[width=\textwidth]{images/fig1_kde_phoenix_v2.png}
    \end{minipage}

    \caption{Densidad KDE de la temperatura máxima diaria para cada
    ciudad y los tres períodos analizados. La línea discontinua vertical
    indica el umbral modal~$\mu^*$ estimado para cada período.}
    \label{fig:kde}
\end{figure}

%% ════════════════════════════════════════════════════════════
\section{Resultados}\label{sec:results}
%% ════════════════════════════════════════════════════════════

El modelo de distribución de mezcla Weibull ajustado mediante
el procedimiento descrito en la Sección~\ref{sec:estimation}
produce un ajuste satisfactorio. El coeficiente de determinación
$R^2$ supera 0.99 en todos los casos. El estadístico de Kolmogorov--Smirnov
(KS) sobre los datos individuales $T_i > 0$ oscila entre 0.015
y 0.048; en la mayoría de los casos el valor~$p$ asociado es
inferior a 0.05, lo que refleja la potencia elevada del test.

\begin{table}[ht]
\centering
\normalem
\begin{tabular}{llrrrrrrr}
\toprule
Ciudad & Período & $\hat\mu^*$ & $k$ & $\lambda$ & $R^2$ & KS & $p_{\mathrm{KS}}$ \\
\midrule
CDMX        & 1940–1968 & 19.54 & 0.884 & 0.367 & 0.9972 & 0.027 & 0.001 \\
        & 1969–1996 & 20.67 & 0.856 & 0.361 & 0.9993 & 0.015 & 0.191 \\
        & 1997–2026 & 22.03 & 0.880 & 0.335 & 0.9971 & 0.023 & 0.006 \\
\midrule
Chihuahua   & 1940–1968 & 30.62 & 0.727 & 0.286 & 0.9978 & 0.035 & 0.005 \\
   & 1969–1996 & 28.50 & 0.730 & 0.172 & 0.9974 & 0.034 & 0.001 \\
   & 1997–2026 & 31.08 & 0.705 & 0.188 & 0.9970 & 0.045 & 0.000 \\
\midrule
Guadalajara & 1940–1968 & 25.00 & 0.862 & 0.244 & 0.9924 & 0.038 & 0.000 \\
 & 1969–1996 & 25.72 & 0.878 & 0.265 & 0.9956 & 0.029 & 0.000 \\
 & 1997–2026 & 27.29 & 0.869 & 0.290 & 0.9961 & 0.029 & 0.000 \\
\midrule
Monterrey   & 1940–1968 & 32.16 & 0.674 & 0.258 & 0.9941 & 0.048 & 0.000 \\
   & 1969–1996 & 29.78 & 0.666 & 0.175 & 0.9951 & 0.048 & 0.000 \\
   & 1997–2026 & 31.95 & 0.652 & 0.196 & 0.9957 & 0.045 & 0.000 \\
\midrule
Phoenix     & 1940–1968 & 38.65 & 0.712 & 0.269 & 0.9961 & 0.045 & 0.001 \\
     & 1969–1996 & 39.09 & 0.699 & 0.237 & 0.9986 & 0.037 & 0.011 \\
     & 1997–2026 & 39.83 & 0.708 & 0.268 & 0.9975 & 0.035 & 0.017 \\
\bottomrule
\end{tabular}
\caption{Parámetros estimados del modelo de mezcla Weibull por ciudad
y período}
\label{tab:parametros}
\end{table}

La comparación entre ciudades revela que Monterrey y Phoenix
presentan los menores valores de $k$ ($0.65$--$0.72$),
lo que indica colas de excedencias relativamente más pesadas:
en estas ciudades, una vez que la temperatura supera el nivel
modal, la probabilidad de excedencias adicionales (\textit{eventos extremos})
decrece más gradualmente. CDMX y Guadalajara muestran los
mayores valores de $k$ ($0.86$--$0.88$), con colas
más ligeras, coherentes con su altitud y clima más templado.

\begin{figure}[ht]
  \centering
  \includegraphics[width=0.5\textwidth]{images/fig6_tendencias_v2.png}
  \caption{Evolución temporal del umbral modal $\mu^*$ por ciudad y
  período. Varias ciudades muestran una tendencia ascendente.}
  \label{fig:tendencias}
\end{figure}

En Ciudad de México, $\mu^*$ aumenta de
$19.54\,{}^\circ\text{C}$ en 1940--1968 a $22.03\,{}^\circ\text{C}$
en 1997--2026, un incremento de $2.49\,{}^\circ\text{C}$ en ocho
décadas. Este es el desplazamiento más pronunciado del conjunto
analizado. Guadalajara muestra un patrón similar, con
un aumento de $25.00$ a $27.29\,{}^\circ\text{C}$
($\Delta\mu^* = 2.29\,{}^\circ\text{C}$). Chihuahua y Monterrey presentan trayectorias
no monótonas: ambas exhiben una reducción en el segundo período
respecto al primero, seguida de una recuperación en el tercero. La tabla \ref{tab:excedencias} presenta las probabilidades estimadas de que la temperatura máxima diaria supere los umbrales
de 30, 35 y 40\,°C.

\begin{table}[ht]
\centering
\normalem
\begin{tabular}{llrrr}
\toprule
Ciudad & Período &
  $P(X>30\,{}^\circ\text{C})$ &
  $P(X>35\,{}^\circ\text{C})$ &
  $P(X>40\,{}^\circ\text{C})$ \\
\midrule
CDMX        & 1940–1968 &  0.28  & 0.016 & 0.001 \\
        & 1969–1996 &  0.37  & 0.015 & 0.001 \\
        & 1997–2026 &  1.50  & 0.109 & 0.007 \\
\midrule
Chihuahua   & 1940–1968 & 27.48 & 2.70  & 0.048 \\
   & 1969–1996 & 25.03  & 3.64  & 0.260 \\
   & 1997–2026 & 38.38  & 8.64  & 0.489 \\
\midrule
Guadalajara & 1940–1968 & 12.26  & 1.75  & 0.211 \\
 & 1969–1996 & 14.12  & 1.98  & 0.235 \\
 & 1997–2026 & 21.29  & 2.54  & 0.239 \\
\midrule
Monterrey   & 1940–1968 & 37.89  & 7.42  & 0.106 \\
   & 1969–1996 & 35.62  & 4.44  & 0.114 \\
   & 1997–2026 & 45.97  & 10.97 & 0.268 \\
\midrule
Phoenix     & 1940–1968 & 47.95 & 32.29 & 19.42 \\
     & 1969–1996 & 48.68 & 33.87 & 14.81 \\
     & 1997–2026 & 51.96 & 36.60 & 17.71 \\
\bottomrule
\end{tabular}
\caption{Probabilidades estimadas de excedencia (\%) de umbrales
de temperatura por ciudad y período.}
\label{tab:excedencias}
\end{table}


Los cambios entre períodos son especialmente notorios para Ciudad de México: $P(X>30\,{}^\circ\text{C})$ aumenta
de 0.28\% en 1940--1968 a 1.50\% en 1997--2026. Este cambio está impulsado principalmente por el desplazamiento de $\mu^*$ de 19.5 a
22.0\,°C.

Para Chihuahua, el aumento en $P(X>35\,{}^\circ\text{C})$
es particularmente pronunciado: de 2.70\% en 1940--1968 a 8.64\% en 1997--2026,
multiplicándose por un factor de 3.2. El aumento en
$P(X>40\,{}^\circ\text{C})$ es aún más pronunciado en términos
relativos: de 0.048\% a 0.489\% (factor 10.2), aunque las
probabilidades absolutas siguen siendo pequeñas.

En Phoenix, un caso naturalmente atípico por su clima árido, las probabilidades de excedencia de umbrales muy altos ($>40\,{}^\circ\text{C}$) presentan un patrón no monótono:
19.4\% (1940--1968), 14.8\% (1969--1996) y 17.7\% (1997--2026).
La reducción en el segundo período puede relacionarse con una
ligera disminución de $k$ (de 0.712 a 0.699) que compensa
parcialmente el aumento de $\mu^*$, aunque también puede reflejar
variabilidad natural en la distribución de los extremos.

\begin{figure}[ht]
    \centering

    % Fila 1
    \begin{minipage}{0.33\textwidth}
        \centering
        \includegraphics[width=\textwidth]{images/fig5_excedencia_cdmx_v2.png}
    \end{minipage}
    \hfill
    \begin{minipage}{0.33\textwidth}
        \centering
        \includegraphics[width=\textwidth]{images/fig5_excedencia_chihuahua_v2.png}
    \end{minipage}
    \hfill
    \begin{minipage}{0.33\textwidth}
        \centering
        \includegraphics[width=\textwidth]{images/fig5_excedencia_phoenix_v2.png}
    \end{minipage}
  \caption{\centering Curvas de probabilidad de excedencia $P(X>\theta)$ por
  ciudad y período}
    \label{fig:excedencia}
\end{figure}

\subsection{Síntesis y contrastación de hipótesis}
\label{subsec:hipotesis}

Los resultados permiten responder a las hipótesis planteadas en
la Sección~\ref{sec:intro}.

\textbf{H1.} Por un lado, se tiene que la forma de la distribución de excedencias no cambia considerablemente en el tiempo dentro de cada
ciudad, pues nótese que el parámetro $k$ varía menos de 3.5\% entre períodos
dentro de cada ciudad. No obstante, la tendencia
decreciente de $k$ en Monterrey (de 0.674 a 0.652)
merece atención: una cola progresivamente más pesada implicaría
que los eventos extremos más intensos se vuelven relativamente
más frecuentes incluso en ausencia de un desplazamiento del nivel
modal, lo que refuerza el riesgo en esa ciudad.

\textbf{H2.}
CDMX, Guadalajara y Phoenix muestran incrementos monotónicos de
$\mu^*$ entre períodos. En estas ciudades, el desplazamiento del
umbral modal es el principal mecanismo para explicar el aumento
en las probabilidades de excedencia. El caso más ilustrativo es
CDMX: una ganancia de 2.5\,°C en $\mu^*$, junto con la estabilidad
de $k$, es suficiente para multiplicar por cinco la
probabilidad de superar 30\,°C.

%% ════════════════════════════════════════════════════════════
\section{Conclusiones}\label{sec:conclusions}

Este trabajo propone y aplica una metodolog\'ia adaptativa para modelar
la cola derecha de la distribuci\'on de temperatura m\'axima diaria
mediante la distribuci\'on Weibull de dos par\'ametros. Los resultados
obtenidos para las cinco ciudades y tres per\'iodos analizados permiten
responder las preguntas planteadas en la introducci\'on.

El umbral modal~$\mu^*$, que representa la temperatura m\'axima m\'as
frecuente, muestra una tendencia ascendente en la mayor\'ia de las
ciudades analizadas. En Ciudad de M\'exico, el incremento acumulado
entre el primer per\'iodo (1940--1968) y el tercero (1997--2026) es de
$2.5\,{}^\circ\mathrm{C}$; en Guadalajara, de
$2.3\,{}^\circ\mathrm{C}$; en Phoenix, de $1.2\,{}^\circ\mathrm{C}$.

Asimismo, el par\'ametro de forma $k$ permanece comparativamente estable dentro
de cada ciudad a lo largo de los tres per\'iodos (con variaciones menores al
4\% en todos los casos), lo que apoya la hip\'otesis~H1. No es posible,
a partir de la metodolog\'ia utilizada, descartar cambios estad\'isticamente
significativos en la forma sin un an\'alisis de inferencia formal
complementario. Los cambios en $\mu^*$ son los m\'as sistem\'aticos
y de mayor magnitud. Los cambios en $\lambda$ son de menor magnitud y
menos monotónicos. El par\'ametro $k$ es el m\'as estable.

Ahora bien, \textit{\textquestiondown Cambian las probabilidades de
superar umbrales cr\'iticos?} S\'i, y en algunos casos de forma
considerable. La probabilidad de superar $30\,{}^\circ\mathrm{C}$ en
CDMX aumenta m\'as de cinco veces entre el primer y el tercer
per\'iodo. En Chihuahua, la probabilidad de
superar $35\,{}^\circ\mathrm{C}$ se triplica entre los mismos
per\'iodos, evidenciando que los climas
semi\'aridos y continentales, con mayor amplitud t\'ermica estival,
presentan una mayor sensibilidad a los desplazamientos del umbral
modal.

Los resultados sugieren que, aunque la forma de la distribuci\'on de
extremos t\'ermicos no cambie sustancialmente, el desplazamiento del
umbral modal hacia temperaturas m\'as altas es suficiente para incrementar
de forma perceptible la probabilidad de superar niveles cr\'iticos. Dado que
la demanda de energ\'ia para climatizaci\'on (especialmente aire
acondicionado) aumenta no linealmente con la temperatura, incrementos
moderados en $\mu^*$ pueden tener consecuencias sobre
la infraestructura energ\'etica urbana. Este efecto es especialmente
relevante en ciudades de r\'apido crecimiento como CDMX, donde la
penetraci\'on del aire acondicionado ha aumentado significativamente en
las \'ultimas d\'ecadas.

%% ════════════════════════════════════════════════════════════
\bibliography{refs_weibull}

\end{document}
